\begin{mistake}{2026-04-12}{07}

\begin{problemcontent}
计算行列式
\[
D=\begin{vmatrix}
0&2&2&2\\
2&0&2&2\\
2&2&0&2\\
2&2&2&0
\end{vmatrix}.
\]

\end{problemcontent}

\begin{solutioncontent}
该行列式各行元素之和都等于 \(6\)，适合先做行变换化简。

对行列式作变换
\[
R_1\leftarrow R_1+R_2+R_3+R_4,
\]
则
\[
D=\begin{vmatrix}
6&6&6&6\\
2&0&2&2\\
2&2&0&2\\
2&2&2&0
\end{vmatrix}.
\]

提出第一行公因子 \(6\)，得
\[
D=6\begin{vmatrix}
1&1&1&1\\
2&0&2&2\\
2&2&0&2\\
2&2&2&0
\end{vmatrix}.
\]

再作行变换
\[
R_2\leftarrow R_2-2R_1,\qquad
R_3\leftarrow R_3-2R_1,\qquad
R_4\leftarrow R_4-2R_1,
\]
得
\[
D=6\begin{vmatrix}
1&1&1&1\\
0&-2&0&0\\
0&0&-2&0\\
0&0&0&-2
\end{vmatrix}.
\]

此时为上三角行列式，故
\[
D=6\cdot 1\cdot(-2)\cdot(-2)\cdot(-2)=-48.
\]

\end{solutioncontent}

\mistakeknowledge{
\begin{itemize}
  \item \textbf{核心公式/定理}：行列式经“某行加上其他行的线性组合”数值不变；上三角行列式等于主对角线元素的乘积。
  \item \textbf{易错点}：提取第一行公因子后，后续行变换要在新矩阵上继续做。不要把“矩阵初等变换”和“行列式值变化规则”混淆。
  \item \textbf{关联知识}：对角元相同、非对角元相同的数值矩阵，常可先利用“各行和相等”构造一行常数，再迅速化为三角形。
\end{itemize}}

\end{mistake}
